\cotefile{70}
\scanpage{17}
% archive sheet number: 16
% source lines 1-9: block cancelled by long diagonal strokes
\[
\deleted{\sigma_i(x_1,\ldots,x_{n+1})=(x_1,\ldots,x_{i-1},\underset{!}{\underline{x_{i+1}}},x_i,\lambda_{i,i+2}x_{i+2},\lambda_{i,i+2}x_{i+3},\ldots,\lambda_{i,i+2}x_{n+1})}
\]
\[
\deleted{H_i:\left\{\begin{aligned}
x_i&=x_{i+1}\\
(\lambda_{i,i+2}-1)x_{i+2}&=\cdots=(\lambda_{i,i+2}-1)x_{n+1}=0
\end{aligned}\right.}
\]
\deleted{Si $H'_i$ est l'\unsure{hyperplan} d'équation $x_i=x_{i+1}$, on a $H_i\subset H'_i$, donc $H_i$ n'est un \unsure{hyperplan} que si $H_i=H'_i$, i.e. les équations $(\lambda_{i,i+2}-1)x_{i+2}=\cdots=(\lambda_{i,i+2}-1)x_{n+1}=0$ sont \unsure{conséquence|conséquences} de $x_i=x_{i+1}$, ce qui bien sûr n'est le cas que si $\lambda_{i,\unsure{i+1|i+2}}=1$.}

% source lines 10-13
\[
\left\{\begin{aligned}
&\lambda_{i,1}(x_{i+2}+\cdots+x_{n+1})=\lambda_{i,2}(x_{i+2}+\cdots+x_{n+1})=\cdots\ \lambda_{i,i-1}(x_{i+2}+\cdots+x_{n+1})=0\\
&\deleted{\unsure{x_{i+1}}}\lambda_{i,i}(x_{i+2}+\cdots+x_{n+1})=x_i-x_{i+1}\\
&\lambda_{i,i+1}(x_{i+2}+\cdots+x_{n+1})=x_{i+1}-x_i\\
&(\lambda_{i,i+2}-1)x_{i+2}=\cdots=(\lambda_{i,i+2}-1)x_{n+1}=0
\end{aligned}\right.\tag*{$(H_i)$}
\]

% source lines 14-16
(en plus, les conditions sur les $\lambda_{i,j}$, $\unsure{0|1}\leq j\leq i+2$, \deleted{\unsure{disant que}} \unsure{exprimant} $\sigma_i^2=\mathrm{id}\ \ldots$)

% source lines 17-19
\deleted{\unsure{Si}} Sur un corps: si on avait $\lambda_{i,i+2}\ne1$, comme $\lambda_{i,i+2}^2=1$, on aurait $\lambda_{i,i+2}=-1\ne1$ (donc car $\ne2$), et les équations de $H_i$ deviennent
\[
\left\{\begin{aligned}
&x_{i+2}=\cdots=x_{n+1}=0\\
&\deleted{\lambda_{i,1}\unsure{\overline{\Sigma}_i|\Sigma_i}(x)=\cdots=\lambda_{i,i-1}\overline{\Sigma}_i(x)=0}\quad\deleted{:\ \overline{\Sigma}_i(x)=x_{i+2}+\cdots+x_{n+1}}\\
&\deleted{\lambda_{i,i}\overline{\Sigma}_i(x)=}\ x_i=x_{i+1}\ \unsure{\}}\\
&\deleted{\lambda_{i,i+1}\overline{\Sigma}_i(x)=x_{i+1}-x_i}
\end{aligned}\right.
\]
% source lines 24-29
ce qui fait $[n+1-(i+1)]+1=n-i+1$ équations \deleted{\unsure{Si} $i\leq n+1$, les premières équations \unsure{étant} $n+1-i$} indépendantes et comme $i\leq n-1$, on trouve une sous-variété linéaire de codim $\geq2$.

Donc on doit avoir $\lambda_{i,i+2}=1$, les équations de $H_i$ deviennent

\scanpage{18}
% archive sheet number: 17
% source lines 1-3
\[
\left\{\begin{aligned}
\lambda_{i,1}\unsure{\overline{\Sigma}_i|\Sigma_i}(x)&=\cdots=\lambda_{i,i-1}\overline{\Sigma}_i(x)=0\\
\lambda_{i,i}\overline{\Sigma}(x)&=x_i-x_{i+1}\\
\lambda_{i,i+1}\overline{\Sigma}(x)&=x_{i+1}-x_i
\end{aligned}\right.
\qquad :\quad\overline{\Sigma}(x)=x_{i+2}+\cdots+x_{n+1}
\]
% source lines 4-5
Si donc les $\lambda_{i,j}$, $1\leq j\leq i-1$, ne sont pas nuls, alors ces équations se réduisent à
\[
\left\{\begin{aligned}
\overline{\Sigma}_i(x)&=0\\
x_i-x_{i+1}&=0
\end{aligned}\right.
\]
% source lines 8-10
ce sont deux équations linéaires indépendantes, donc $H_i$ serait de codim $2$, donc pas un hyperplan. \unsure{On veut|On doit} donc
\[
\lambda_{i,1}=\cdots=\lambda_{i,i-1}=0
\]
% source lines 12-13
et les équations se \deleted{\unsure{se}} réduisent alors:
\[
\lambda_{i,i}\overline{\Sigma}_i(x)=-\lambda_{i,i+1}\overline{\Sigma}_i(x)=x_i-x_{i+1}
\]
% source lines 14-18
\deleted{Si on a} \deleted{\unsure{On avait} $\lambda_{i,i}=\lambda_{i,i+1}=0$, ces équations}

Les relations exprimant $\sigma_i^2=\mathrm{id}$ se réduisent ici à: $\lambda_{i,i}+\lambda_{i,i+1}=0$, de sorte qu'il reste un seul paramètre $\lambda_{i,i+1}=1+\alpha_i$, et on trouve
% source bottom matrix
\[
\left\{\begin{aligned}
&\deleted{\sigma_i=}\\[-1.5ex]
\sigma_i={}&
\begin{array}{c@{\,}l}
\begin{array}{r}
\left.\vphantom{\begin{matrix}1\\1\\1\end{matrix}}\right\}\unsure{i-1}\\
\left.\vphantom{\begin{matrix}1\\1\end{matrix}}\right\}2\\
\left.\vphantom{\begin{matrix}1\\1\\1\\1\end{matrix}}\right\}\deleted{\unsure{n-i}}\\
1
\end{array}
&\left(\begin{array}{cccccccccc}
1&&0&&&&&&&0\\
&1&&&&&0&&&0\\
0&&\ddots&&&&&&&\vdots\\
&&&1&&&&&&\\
&&&&\multicolumn{2}{c}{\left(\begin{matrix}0&1\\1&0\end{matrix}\right)}&\begin{matrix}-(1+\alpha_i)\deleted{\unsure{x}}\\(1+\alpha_i)\deleted{\unsure{x}}\end{matrix}&\begin{matrix}-(1+\alpha_i)\\(1+\alpha_i)\end{matrix}&\begin{matrix}\cdots\;-(1+\alpha_i)\\\cdots\;(1+\alpha_i)\end{matrix}&\begin{matrix}0\\0\end{matrix}\\
&&&&&&\deleted{\unsure{1}}&&&\deleted{\unsure{0}}\\
&&&&&&1&&0&0\\
&&&&&&&1&\deleted{\unsure{\cdots}}&\unsure{0}\\
&&&&&&&&\ddots\ \deleted{\unsure{0}}&\vdots\\
&&&&&&&&1&0\\
0&0&\cdots&0&0&0&0&\cdots&0&1
\end{array}\right)
\end{array}
\end{aligned}\right.
\]

\scanpage{19}
% archive sheet number: 18
% source top cancellation and lines 1-3
\deleted{$\unsure{\sigma}$}
\[
\left\{\begin{aligned}
\sigma_i(x_1,\ldots,x_{n+1})={}&(x_1,\ldots,x_{i-1},x_{i+1}-(1+\alpha_i)(x_{i+2}+\cdots+x_{n+1}),x_i+(1+\alpha_i)(x_{i+2}+\cdots+x_{\unsure{n+2|n+1}}),\\
&\hspace{6em}x_{i+2},\ldots,x_{n+1})\\[2ex]
H_i:\quad &x_i-x_{i+1}+(1+\alpha_i)(x_{i+2}+\cdots+x_{\unsure{n+2|n+1}})=0
\end{aligned}\right.
\]
% source lines 4-7
Ces résultats marchent aussi sur tt anneau \unsure{réduit}. Sur un anneau quelconque [il y a lieu de \deleted{faire} poser
\[
\lambda_{i,\unsure{i+2|i+1}}=1+\unsure{\nu|v}\qquad \nu\ \text{nilpotent}
\]
% source lines 8-12
\deleted{et les relations} en car \unsure{rés.}\ $\ne2$ les relations qui expriment $\sigma_i^2=\mathrm{id}$ deviennent \unsure{$\nu$}
\[
\left\{\begin{aligned}
&\inserted{\nu=0}\\[-.5ex]
&\lambda_{i,1}=\cdots=\lambda_{i,i-1}=0\\
&\lambda_{i,i}\deleted{\unsure{\nu}}+\lambda_{i,i+1}\deleted{+\nu\lambda_{i,i}}=\lambda_{i,i}+\lambda_{i,i+1}\deleted{+\nu\lambda_{i,i+1}}=0\\
&\deleted{\unsure{\nu^2+\nu}}
\end{aligned}\right.
\]
\deleted{donc $\lambda_{i,i}+\lambda_{i,i+1}=-\nu\lambda_{i,i}=-\nu(\lambda_{i,i+1})$}

% source lines 13-16
Comme devant. \unsure{Il reste: examinons} ce qui se passe en car rés. $2$ (sur un anneau entier, \unsure{$2$ non unité}), où on trouve, \deleted{\unsure{et}}
\[
\left\{\begin{aligned}
&\underset{\text{\unsure{nilpotent|inversible}}}{\underline{(2+\nu)}}\lambda_{i,1}=(2+\nu)\lambda_{i,2}\ \cdots=(2+\nu)\lambda_{i,i-1}=0\\
&\lambda_{i,i}+\lambda_{i,i+1}=-\nu\lambda_{i,i}=-\nu\lambda_{i,i+1}\quad\deleted{\unsure{\nu\lambda_i}}\\
&\deleted{\unsure{\nu}}\ \nu^2+2\nu=0\qquad\text{i.e.}\quad\nu(2+\nu)=0
\end{aligned}\right.
\]
% source bottom consequences; arrow from the system to the left equation
\[
\begin{aligned}
\swarrow\quad&\\[-1ex]
\lambda_{i,i+1}&=-(1+\nu)\lambda_{i,i}\quad\deleted{\unsure{-\nu\lambda_{i,i}=}}\qquad
\nu\underbrace{[\lambda_{i,i+1}-\lambda_i]}_{-(2+\nu)\lambda_{i,i}\ \deleted{\unsure{\nu}}}=0
\end{aligned}
\]

\scanpage{20}
% archive sheet number: 19
% source top system, lines 1-3
\[
\unsure{\Downarrow}\quad
\left\{\begin{aligned}
(2+\unsure{\nu|v})\lambda_{i,1}&=\cdots=(2+\nu)\lambda_{i,i-1}=0\\
\lambda_{i,i+1}&=-(1+\nu)\lambda_{i,i}\\
\nu(2+\nu)&=0
\end{aligned}\right.\quad\biggr]
\]
% source lines 4-6
Reprenons les équations de $H_i$. $H_i$ est contenu dans l'hyperplan $H'_i$ d'équation
\[
\lambda_{i,i}\unsure{\overline{\Sigma}_i|\Sigma_i}(x)=x_i-x_{i+1}
\]
% source lines 7-11
Si \deleted{on avait} $H_i$ était un \unsure{hyperpl.}, on aurait $H_i=H'_i$, i.e. les autres équations de $H_i$ seraient conséquences de la précédente. Ce n'est possible manifestement que si $\lambda_{i,1}=\cdots=\lambda_{i,i-1}=0$, $\lambda_{i,i}+\lambda_{i,i+1}=0$, $\lambda_{i,i+2}=1$, \unsure{OK}\ $\ldots$

% source horizontal separator
\noindent\rule{5cm}{.4pt}\par\noindent\rule{4cm}{.4pt}

% source lines 12-15
\uline{NB} $\sigma_0$ ne change que la première coordonnée, $\sigma_i$ ($1\leq i\leq\unsure{n-1|n}$) que les $i$ème et $(i+1)$ème. Ceci \deleted{\unsure{montre}}\inserted{implique que} $\sigma_i$ et $\sigma_j$ commutent si $j\geq i+2$

% source lines 16-19: right ends reach the paper edge
\[
\begin{aligned}
\sigma_0\sigma_1(\deleted{\unsure{x}}x_1,x_2,x_3,\deleted{\unsure{x}}\ldots,x_{n+1})
&=\sigma_0(x_2-(1+\alpha_1)\Sigma_1,x_1+(1+\alpha_1)\Sigma_1,x_3,\unsure{\ldots}\\
&=(1-x_2+(1+\alpha_1)\Sigma_1+\alpha_0\unsure{(}x_1+(1+\alpha_1)\Sigma_1+\Sigma_1),x_1+(1+\alpha_1)\unsure{\Sigma_1},\\
&\hspace{4em}x_3,\ldots,x_{n+1})\\
&\deleted{\unsure{\ne}}\\[-1.2ex]
&=(1+\alpha_0x_1-x_2+\deleted{\unsure{[(1+\alpha_0)(1+\alpha_1)+\alpha_0]}}\inserted{(1+\alpha_1+2\alpha_0+\alpha_0\alpha_1)}\Sigma_1,x_1+(1+\alpha_1)\Sigma_1,x_3,\unsure{\ldots,x_{n+1})}
\end{aligned}
\]
% source lines 20-21
\[
\begin{aligned}
(\sigma_0\sigma_1)^p(x_1,x_2,\ldots,x_{n+1})={}&(F_p(\alpha_0,\alpha_1;x_1,\ldots,x_{n+1}),G_p(\alpha_0,\alpha_1;\unsure{x_1|\alpha_0},\ldots,\unsure{x_{n+1}}),\\
&\hspace{6em}x_3,\ldots,x_{n+1})
\end{aligned}
\]
% source bottom line
Donc $(\sigma_0\sigma_1)^{\unsure{p}}\inserted{=\mathrm{id}}\Longleftrightarrow$

\scanpage{21}
% archive sheet number: 20
% source lines 1-3
\[
(i\geq1)\quad\left\{\begin{aligned}
\sigma_i\mid\unsure{f_{i-1}|\mathfrak f_{i-1}}&=\mathrm{id}\\
\sigma_i\mid(f_{i+1}/f_{i-1})&=\text{échange des coordonnées}\quad x_i,x_{i+1}\\
\sigma_i\mid E/f_{i+1}&=\mathrm{id}
\end{aligned}\right.
\]
% source lines 4-5
\[
\left\{\begin{aligned}
\sigma_0\mid f_1&=\text{échange de}\quad\unsure{b_0|s_0},b_1\\
\sigma_0\mid E/f_1&=\mathrm{id}
\end{aligned}\right.
\]
% source left margin beside lines 6-8
% slanted
\margin{\deleted{\unsure{$i\geq0$}}}
% source lines 6-8
\[
(i\geq1)\quad\left\{\begin{aligned}
\sigma_i\sigma_{i+1}\mid f_{i-1}&=\mathrm{id}\\
\sigma_i\sigma_{i+1}\mid f_{i+2}/f_{i-1}&\overset{?}{=}\\
\sigma_i\sigma_{i+1}\mid E/f_{i+2}&=\mathrm{id}
\end{aligned}\right.
\]
% source lines 9-10
\[
\left\{\begin{aligned}
\sigma_0\sigma_1\mid f_2&\overset{?}{=}\\
\sigma_0\sigma_1\mid E/f_2&=\mathrm{id}
\end{aligned}\right.
\]

\scanpage{22}
% archive sheet number: 21
% source lines 1-7
Soit $G$ l'image de $\unsure{G_n|\mathscr G_n}$ dans $\mathrm{Aut.\,Aff}(E)$ (on suppose maintenant que $k$ est un corps), $\Pi$ la \deleted{polyèdre}\inserted{$n$-carte} \uline{régulière} \deleted{\unsure{primitive}} \unsure{correspondante|correspondant}, On a donc une représentation \uline{régulière} de $\Pi$. \unsure{Correspondant|Correspondent}: une représentation \uline{fidèle} de $G=\mathrm{Aut}(\Pi)$ dans $\mathrm{Aut\,Aff}(E)$, \deleted{\unsure{et}} une application \uline{injective} de $\unsure{\mathrm{Drap}(\Pi)|\mathrm{Rep}(\Pi)}$ dans $\unsure{\mathrm{Drap.\,max}(E)|\mathrm{Drap}_{\max}(E)}$.

% source lines 8-10
Les applications $D_i(\Pi)\longrightarrow\deleted{\unsure{\mathrm{Esp.}}}\unsure{\mathrm{Grass.aff}_i(E)|\mathrm{Grass}_i(E)}$ sont elles aussi injectives -- en d'autres termes, si \deleted{$\unsure{\sigma,\sigma'}\in$} $g\deleted{,\unsure{g'}}\in G\subset\mathrm{Aut.\,Aff}(E)$
% source lines 11-12: cancelled wording and its interlinear addition
\deleted{(\unsure{le|G} \unsure{engendré par} $\sigma_0,\ldots,\sigma_n$) \unsure{est tel}} que \deleted{\unsure{g et g' coïncident} \inserted{\unsure{fixant le drapeau}} \unsure{en un même} $g\cdot\unsure{f_0|\mathfrak f_0},\unsure{f_1},\ldots,\unsure{f_n}$}
% source lines 13-20
$gf_i=f_i$, $g$ est-il dans le groupe engendré par $\sigma_0\,\cdots\,\sigma_{i-1}\,\widehat{\sigma_i}\,\sigma_{i+1}\,\cdots\,\sigma_n$? Ce n'est pas évident du tout -- et c'est pourquoi il n'est pas clair que des résultats sur les \deleted{\unsure{ordres}} ordres des polygones associés à la \unsure{réal.} géom \deleted{de}\inserted{\unsure{associée|envisagée}} $\Pi$, impliquent quelque chose sur les ordres des $\sigma_i\sigma_{i+1}$!

\scanpage{23}
% archive sheet number: 22
% source lines 1-2
\[
\underbrace{\sigma_0\sigma_1}_{\unsure{\pi|\Pi}}(x_1,\unsure{x_2},x_3,\ldots,x_n)
=(1+\alpha_0x_1-x_2+\underset{\substack{\Vert\\1+\alpha_1+2\alpha_0+\alpha_0\alpha_1}}{a}\Sigma_1,
x_1+\underset{\substack{\Vert\\1+\alpha_1}}{b}\Sigma_1,x_3,\ldots)
\]
% source line 3, cancelled
\[
\deleted{\unsure{\pi^{-1}}=\sigma_1\sigma_0(x_1,\unsure{x_2},\ldots,\unsure{x_n})=\unsure{\sigma|0}}
\]
% source lines 4-7
\[
\begin{aligned}
\pi^2(x_1,\ldots,x_n)={}&\underline{1}+\alpha_0(\underline{1}+\underline{\alpha_0x_1}-\underline{x_2}+\underline{a\Sigma_1})-\underline{x_1}-\underline{b\Sigma_1}+\underline{a\Sigma_1},\\
&1+\alpha_0x_1-x_2+a\Sigma_1+b\Sigma_1,x_3,\ldots)\\
={}&((1+\alpha_0)+(\alpha_0^2-1)x_1-\alpha_0x_2+(a\alpha_0+a-b)\Sigma_1,\\
&1+\alpha_0x_1-x_2+(a+b)\Sigma_1,x_3,\ldots)
\end{aligned}
\]
% source lines 8-10
\[
\begin{aligned}
\pi^3(x_1,\ldots,x_n)={}&1+\alpha_0(1+\alpha_0)+\alpha_0(\alpha_0^2-1)x_1-\alpha_0^2x_2+\alpha_0(a\alpha_0+a-b)\Sigma_1\\
&-1\deleted{\unsure{+\alpha}}-\alpha_0x_1+x_2-(a+b)\Sigma_1+a\unsure{\Sigma_1},\\
&\underbrace{(1+\alpha_0)+(\alpha_0^2-1)x_1-\alpha_0x_2+\unsure{[}a(\alpha_0+1)\Sigma_1}_{\phantom{(1+\alpha_0)+(\alpha_0^2-1)x_1}} ,x_3,\unsure{\ldots}
\end{aligned}
\]
% source lines 11-12; curved forwarding arrow reproduced at right
\diagram{\begin{tikzcd}[ampersand replacement=\&,column sep=small,row sep=small]
{}\&{}\arrow[dd,bend left=20]\&{}\\
{\begin{aligned}
=\deleted{\unsure{1+}}\ \alpha_0(1+\alpha_0)+\alpha_0(\alpha_0^2-2)x_1+(1-\alpha_0^2)x_2
+\unsure{(}\deleted{\unsure{a\alpha_0^2+a\alpha_0}}\qquad\qquad\\
(\alpha_0+1)(a\alpha_0-b)\deleted{\unsure{+b-a|+\alpha_0-1}})\Sigma_1,
\end{aligned}}\&{}\&{}\\
{}\&{}\&{}
\end{tikzcd}}
% source left note alongside final calculation
\margin{$\alpha_0=-1$}
% source lines 13-14
\[
\left\{\begin{aligned}
\alpha_0(1+\alpha_0)&=0\qquad\alpha_0(\alpha_0^2-2)=1\qquad\alpha_0^2=1\quad\bigm|\quad\alpha_0=-1\\[2ex]
&\deleted{\unsure{2a}}
\end{aligned}\right.
\]

\scanpage{24}
% archive sheet number: 23
% source lines 1-4, cancelled trials
\[
\deleted{u(x,y)=\sigma_0\sigma_1(x,y)=}
\]
\[
\deleted{\left\{\begin{aligned}
\sigma_0(x,y)&=(1-(1+\unsure{\nu_0|v_0})x+\alpha_0y,y)\\
\sigma_1(x,y)&=(-\nu_1y,x+(\nu_1+1)y)
\end{aligned}\right.}
\]
\[
\deleted{\unsure{u(x,y)=\sigma_0\sigma_1(x,y)=}}
\]
% source lines 5-7
\[
\sigma_0(x_1,\ldots,x_{n+1})=(1-x_1+\alpha_0\Sigma_0,x_2,\ldots,x_{n+1})
\]
\[
\sigma_i(x_1,\ldots,x_{n+1})=(x_1,\ldots,x_{i-1},x_{i+1}-(1+\alpha_i)\Sigma_i,x_i+(1+\alpha_i)\Sigma_i,x_{\unsure{i+3|i+2}},\ldots,x_{n+1})\tag*{$1\leq i\leq n-1$}
\]
\[
\sigma_n(x_1,\ldots,x_{n+1})=(x_1,\ldots,x_{n-1},x_{n+1},x_n)
\]
% source line 8, cancelled quadratic expression
\[
\deleted{f(x_1,\ldots,x_{n+1})=\sum_{\unsure{0|1}\leq i\leq n}a_{\deleted{\unsure{j}}i}^2x_i^2\unsure{a_j|x_j}+\sum_{1\leq i<j\leq\unsure{n+1}}\unsure{d_{ij}|a_{ij}}x_ix_j+\sum_{1\leq i\leq n+1}c_ix_i}
\]
% source lines 9-14: boxed block cancelled with diagonal strokes
\deleted{\uline{\unsure{Incidence|Invariants}?}}
\[
\begin{array}{|l}\hline
\deleted{\begin{aligned}
H_0:\quad&2x_1\unsure{-|+}\alpha_0(x_2+\cdots+x_{n+1})=z\\
D_0:\quad&x_2=x_3=\cdots=x_{n+1}=z=0\\[2ex]
H_i:\quad&x_i-x_{i+1}+(1+\alpha_i)(x_{i+2}\deleted{\unsure{=0}}+x_{i+3}+\cdots+x_{n+1}\unsure{\}|)}=0\\
D_i:\quad&\{x_1=\cdots\ x_{i-1}=x_{i+2}=\cdots=x_{n+1}=z=0
\end{aligned}}\\\hline
\end{array}
\]
\margin{\deleted{(pt à l'infini de \deleted{la droite} $\unsure{\mathrm{dr}(s_0s_1)|\mathrm{dr}(b_0b_1)|\mathrm{dr}(s_0,s_1)}$\unsure{)}}}
% source lines 15-16
\[
\begin{aligned}
\overline{\sigma}_0(\underbrace{x_1,\ldots,x_{n+1},\overset{x_0}{\overset{=}{z}}}_{\overline{x}})
&=(x_0-x_1+\alpha_0(x_2+\cdots+x_{n+1}),x_2,\ldots,x_{n+1},x_0)\\
&=\overline{x}+\deleted{\unsure{e_1}}[e'_0-e'_1+\alpha_0(e'_2+\cdots+e'_{n+1})]\unsure{\otimes}e_1
\end{aligned}
\]
% source line 17
\[
\overline{\sigma}_i(\overline{x})=\overline{x}+[e'_i-e'_{i+1}+(1+\alpha_i)(e'_{i+2}+\cdots+e'_{n+1})]\otimes(-e_i+e_{i+1})\tag*{$(1\leq i\leq n)$}
\]
% source line 18
\[
\overline{\sigma}_n(\overline{x})=\overline{x}+(e'_n-e'_{n+1})\otimes(-e_n+e_{n+1})
\]

\scanpage{25}
% archive sheet number: 24
% source top interlinear note and line 1
\[
H_0:\quad\overset{\substack{z\ \text{variable d'\unsure{homogénéité}}\\\Vert}}{x_0}-2x_1+\alpha_0(x_2+\cdots+x_{n+1})=0
\]
% source line 2
$\deleted{D_0}\unsure{\delta_0|d_0}:$ pt à l'infini dans la direction de $e_1=s_1-s_0$

% source line 3
\[
H_i:\quad x_i-x_{i+1}+(1+\alpha_i)(x_{i+2}+\cdots+x_{n+1})=0
\]
% source left of lines 3-5
\margin{$1\leq i\leq n$}
% source lines 4-5
$\delta_i:$ pt à l'infini dans la direction de
\[
\begin{aligned}
e_{i+1}-\unsure{e_i}&\\[-1ex]
&=s_{i+1}-s_i
\end{aligned}
\]
\noindent\rule{3cm}{.4pt}

\scanpage{26}
% archive sheet number: 25
% source landscape, top lines 1-3
\[
\sigma_0\unsure{s_i}=-(1+\unsure{\alpha_i|\alpha_0})s_0+(1+\alpha_i)s_1+s_i
\]
\[
i\geq2
\]
\[
\sigma_0s_{i+1}=\sigma_0\sigma_is_i=\sigma_i\sigma_0s_i=-(1+\alpha_i)s_0+(1+\alpha_i)s_1+s_{i+1}
\]
% source middle, lines 4-6
\[
\begin{aligned}
\sigma_is_j&=s_j\qquad 0\leq j\leq i-1\\
\sigma_is_i&=s_{i+1}\\
\sigma_is_{i+1}&=s_i
\end{aligned}
\]
% source middle, lines 7-8
\[
\begin{aligned}
\sigma_is_{i+2}&=(\underbrace{0,\ldots,0}_{i-1},-(1+\alpha_i),(1+\alpha_i),1,0,0,\ldots)=s_{i+2}+(1+\alpha_i)(s_{i+1}-s_i)\\
&=-(1+\alpha_i)s_i+(1+\alpha_i)s_{i+1}+s_{i+2}
\end{aligned}
\]
% source lower box heading
\deleted{$\unsure{\sigma_j|\sigma_i}s_{i+2}$}
\fbox{Pour $0\leq i\leq n$}
% source lower boxed case formula
\[
\boxed{\sigma_is_j=\left\{\begin{aligned}
&s_j&&\text{si }0\leq j\leq i-1\\
&s_{i+1}&&\text{si }j=i\\
&s_i&&\text{si }j=i+1\\
&-(1+\alpha_i)s_i+(1+\alpha_i)s_{i+1}+s_j&&\text{si }i+2<j\leq n+1
\end{aligned}\right.}
\]
% source bottom left box, lines 13-14
\deleted{$\sigma_is_j-s_j=$}
\[
\boxed{\begin{aligned}
\sigma_i-\mathrm{id}&=\unsure{h_i|l_i}\otimes(s_{i+1}-s_i)\qquad\deleted{\unsure{\sigma_i s_j=x-x}\{}\\
h_i&=s'_i-s'_{i+1}+(1+\alpha_i)(s'_{i+2}+\cdots+s'_{n+1})
\end{aligned}}
\]

\[
\sigma_is_j-s_j=\left\{\begin{aligned}
&0&&\text{si }\unsure{i|0}\leq j\leq i-1\\
&s_{i+1}-s_i&&\text{si }j=i\\
&-(s_{i+1}-s_i)&&\text{si }j=i+1\\
&(1+\alpha_i)(s_{i+1}-s_i)&&\text{si }j\geq i+2
\end{aligned}\right.
\]

\scanpage{27}
% archive sheet number: 26
% source lines 1-7
Pour modifier cette situation et ouvrir la voie à une politique vraiment nouvelle il faut affaiblir de manière décisive l'emprise politique et économique des monopoles. Toute autre perspective serait illusoire. En refusant les nationalisations indispensables, en proposant de laisser le capital privé maître de ses filiales, en réduisant les mesures fiscales contre les milliardaires, le PS et le MRG vident le programme commun de son efficacité.

% source lines 8-14
Nous partageons la déception et l'inquiétude de tous ceux qui ont placé leur espoir dans le programme commun. Cependant une voie existe pour assurer la victoire de la gauche, non pas seulement la victoire électorale mais le succès de la politique anti-monopoliste à laquelle elle ne peut renoncer sans se désavouer. La gauche française ne doit pas se rabattre sur une politique de gestion de la crise, à la manière des partis socio-démocrates au pouvoir.

% source lines 15-17
La seule voie est la lutte immédiate pour obtenir la signature d'un programme commun actualisé, c'est-à-dire ajusté à une situation très aggravée depuis 1972.

% source lines 18-21
Nous appelons tous les membres du personnel des Universités à prendre leurs responsabilités. Dans ce moment crucial l'avenir dépend de vous. Plus que jamais c'est vous, c'est nous qui pouvons faire changer le Parti Socialiste.

% source lines 22-24
C'est pourquoi nous vous appelons à soutenir nos propositions non parce que ce sont celles du Parti Communiste mais parce qu'elles sont désormais les seules fidèles au programme de 1972 et à sa portée transformatrice.

% source line 25
\begin{center}
POUR LE CHANGEMENT PAR L'UNION SUR UN BON PROGRAMME
\end{center}

% source lines 26-27
\begin{center}
MANIFESTEZ LE 10 DECEMBRE 1977 A 15 H à MONTPELLIER
\end{center}
\begin{flushright}
(rassemblement aux Arceaux)
\end{flushright}

% source line 28
\begin{flushright}
Section des Universités du Parti Communiste Français
\end{flushright}

\scanpage{29}
% archive sheet number: 28
% source upper centre
(8 p)

% source upper right, two title lines
\begin{flushright}
$n$-quasi-polyèdres

(\unsure{vieille rédaction|vieilles rédactions})
\end{flushright}

\scanpage{30}
% archive sheet number: 29
% source top heading
\uline{$n$-Quasi-polyèdres réguliers en coordonnées can.}

% source lines 1-2
\[
\left\{\begin{aligned}
&\sigma_0,\ldots,\sigma_n\in\mathrm{Aff}(E)\qquad E\ \text{affine de rang }n+1\\
&\sigma_j\sigma_i=\sigma_i\sigma_j\quad\text{si }0\leq i\leq i+2\leq j\leq n
\end{aligned}\right.\tag*{(1)}
\]
% source line 3
\[
E_0=\{s_0\}\subset E_1\subset\cdots\subset E_n\qquad\text{drapeau de }E\tag*{(\unsure{1|2})}
\]
% source lines 4-5
\[
\left\{\begin{aligned}
\sigma_iE_j&=E_j\quad\text{si }j\ne i\\
\sigma_iE_i&\ne E_i
\end{aligned}\right.\tag*{(3)}
\]
% source lines 6-10
\[
\left\{\begin{aligned}
&s_0\\
&s_1=\sigma_0s_0\\
&\cdots\\
&s_i=\sigma_{i-1}s_{i-1}\qquad(1\leq i\leq n+1)\\
&\cdots\\
&s_{n+1}=\sigma_ns_n
\end{aligned}\right.\tag*{(4)}
\]
% source lines 11-14
\uline{Prop} $(s_0,\ldots,s_{n+1})$ est une base affine de $E$, et
\[
E_i=\unsure{\mathrm{Env.aff}|\mathrm{Ens.aff}}(s_0,\ldots,s_i)\tag*{(5)}
\]
\[
\sigma_ks_i=s_i\quad\text{si }0\leq i<k\leq n\deleted{\unsure{+1}}\tag*{(6)}
\]
\inserted{\uline{pour $\unsure{i}$, $0\leq i\leq n+1$}}

% source lines 15-17
Par induction \unsure{(on va prouver} que $(s_0,\ldots,s_i)$ est aff\textsuperscript{t} libre et (5), (6). Trivial si $i=0$. Supposons \inserted{$i\geq1$ et} prouvé pour $j<i$. Il faut prouver
% source lines 18-20
\[
\left\{\begin{aligned}
&s_i\notin E_{i-1}(=\mathrm{Env}(s_0,\ldots,s_{i-1})\ \text{par hyp. réc})\\
&E_i=\mathrm{Env}(E_{i-1},s_i)
\end{aligned}\right.
\]
\inserted{$\sigma_{\unsure{h|k}}s_i=s_i$ pr $i<h\leq n$}
% source lines 21-23
Si on avait $s_i\in E_{i-1}$, i.e. $\sigma_{i-1}s_{i-1}\in E_{i-1}$, comme $\sigma_{i-1}s_j=s_j$ pour $0\leq j\leq i-2$ et $E_{i-1}=\mathrm{Env}(s_0,\ldots,s_{i-1})$ par réc, on aurait
\[
\sigma_{i-1}E_{i-1}=\mathrm{Env}(\underset{\substack{\Vert\\s_0}}{\sigma_{i-1}s_0},\ldots,\underset{\substack{\Vert\\s_{i-2}}}{\sigma_{i-1}s_{i-2}},\underset{\substack{\uparrow\\E_{i-1}}}{\sigma_{i-1}s_{i-1}})\subset E_{i-1}
\]
% source lines 25-28
donc $\sigma_{i-1}E_{i-1}=E_{i-1}$, contrairement à (3 b), OK.

On a, comme $\sigma_{i-1}E_i=E_i$ et $s_{i-1}\in E_{i-1}$ (réc) $\subset E_i$, $\sigma_{i-1}(s_{i-1})\subset E_i$ i.e. $s_i\in E_i$. Comme $s_i\notin E_{i-1}$, on a donc $E_i=\mathrm{Env}(E_{i-1},s_i)$. Enfin, on a pour \deleted{$h>i$}\inserted{$i<h\leq n$}

\[
\sigma_hs_i=\sigma_h\sigma_{i-1}s_{i-1}=\sigma_{i-1}(\sigma_hs_{i-1})\overset{\text{réc}}{=}\sigma_{i-1}(s_{i-1})=s_i
\]
OK.

\scanpage{31}
% archive sheet number: 30
% source lines 1-4
On utilise les $s_i$ ($0\leq i\leq n+1$) comme base de $E$, \unsure{de|donc} $s_0$ comme origine, et les
\[
e_i=s_i-s_0\qquad1\leq i\leq n+1\tag*{(7)}
\]
comme base de l'espace des \unsure{translations} $V$ de $E$.

% source lines 5-8
On veut expliciter les matrices (à $n+2$ lignes et $n+2$ colonnes) des $\sigma_i$ $0\leq i\leq n$. Revient à la connaissance des $\sigma_is_j$ $0\leq i\leq n$, $0\leq j\leq n+1$. Je dis qu'il suffit de connaître
% source lines 9-10
\[
\underset{\substack{\Vert\\\sigma_i^2s_i}}{\sigma_is_{i+1}}=\inserted{u_i},\quad\sigma_is_{i+2}=\unsure{v_i|\nu_i}\tag*{(8)}
\]
% source lines 11-15
Car on a \deleted{\unsure{a}}
\[
\left\{\begin{aligned}
&\sigma_is_j=s_j\quad\text{si }0\leq j<i\\
&\sigma_is_i=s_{i+1}\\
&\sigma_is_{i+1}=u_i\in E_{i+1}\qquad E_{i+1}=\mathrm{Env}(E_{i-1},s_{i+1},u_i)\\
&\sigma_is_{i+2}=v_i\in E_{i+2}\qquad E_{i+2}=\mathrm{Env}(E_{i+1},v_i)\quad\inserted{(0\leq i\leq n-1)}\\
&\sigma_is_j=\sigma_{j-1}\cdots\sigma_{i+2}(v_i)\quad\deleted{\unsure{k}}\ \text{si }\unsure{i+3}\leq j\leq n+1
\end{aligned}\right.\tag*{(9)}
\]
% source lines 16-17
Car $s_j=\sigma_{j-1}\cdots\sigma_{i+2}s_{i+2}$ donc
\[
\sigma_is_j=\sigma_i\sigma_{j-1}\cdots\sigma_{i+2}s_{i+2}=\sigma_{j-1}\cdots\sigma_{i+2}\underset{v_i}{\underline{\sigma_is_{i+2}}}
\]
% source lines 18-20
Posons
\[
\deleted{\sigma_is_{i+1}=s_0+}
\]
\[
\left\{\begin{aligned}
u_i&=(u_{i,1},\ldots,u_{i,i+1})=s_0+\sum_{1\leq\deleted{\unsure{k}}h\leq i+1}u_{i,h}e_h\\
v_i&=(v_{i,1},\ldots,v_{i,i+2})=s_0+\sum_{1\leq h\leq i+2}v_{i,h}e_h
\end{aligned}\right.\tag*{(10)}
\]
% source lines 21-23
de sorte qu'on obtient si $j\geq i+3$
\[
\begin{aligned}
\sigma_is_j&=\sigma_{j-1}\cdots\sigma_{i+2}(s_0+v_{i,1}e_1+\cdots+v_{i,i+1}e_{i+1}+v_{i,i+2}e_{i+2})\\
&=s_0+\sum_{0\leq h\leq i+1}v_{i,\unsure{k|h}}e_h+v_{i,i+2}e_j
\end{aligned}
\]
% source bottom, lines 24-26
i.e.
\[
\left\{\begin{aligned}
\sigma_is_j&=(v_{i,1},\ldots,v_{i,i+1},0,\ldots,0,\overset{j}{v_{i,i+2}},0,\ldots,0)\\
&=s_0+\sum_{0\leq h\leq i+1}v_{i,h}e_h+v_{i,i+2}e_j\deleted{\unsure{n}}\qquad\text{si }i+2\leq j\leq n+1
\end{aligned}\right.\tag*{(10)}
\]

\scanpage{32}
% archive sheet number: 31
% source lines 1-2
compte tenu que, pr $0\leq i\leq n\deleted{\unsure{+1}}$, $\sigma_i e_i=e_{i+1}$. On a plus explicitement
% source lines 3-11
\[
\left\{\begin{aligned}
&\sigma_i e_j=e_j\quad\text{si }1\leq j<i\\
&\sigma_i e_i=e_{i+1}\ \deleted{\unsure{s}}\quad(\text{si }i\geq1)\deleted{\unsure{\leq n}}\\
&\sigma_i e_{i+1}=\deleted{\unsure{u_i-s_0}}\left\{\begin{aligned}
&u_0-s_1=(u_0-s_0)\deleted{\unsure{+}}-e_1=(u_{01}-1)e_1&&\text{si }i=0\\
&u_i-s_0=\sum_{1\leq h\leq i+1}\deleted{\unsure{u_{i,1}}}u_{i,h}e_h&&\text{si }i\geq1
\end{aligned}\right.\\
&\deleted{\unsure{\sigma_i e_{i+2}}=\{\unsure{v_0|\nu_0}-s_1}\\
&\left[\sigma_0 e_{i+2}=\left\{\begin{aligned}
&v_0-s_1=(v_0-s_0)-e_1=(v_{01}-1)e_1+v_{02}e_2&&\text{si }i=0\\
&\unsure{v_0|v_i}-s_0=\sum_{1\leq h\leq i+2}v_{i,h}e_h&&\text{si }i\geq1
\end{aligned}\right.\right]\\
&\sigma_i e_j=\left.\left\{\begin{aligned}
&(v_{01}-1)e_1+v_{02}e_j\ \deleted{\unsure{+}}&&\text{si }i=0\\
&\deleted{\unsure{v_{i,1}}}\sum_{1\leq h\leq i+1}v_{i,h}e_h+v_{i,i+2}e_j&&\text{si }i\geq1
\end{aligned}\right.\right\}\ \deleted{\unsure{j}}\ \inserted{\text{si }i+2\leq j\leq n+1}
\end{aligned}\right.\tag*{(11)}
\]

% source lines 12-13
\deleted{\unsure{Commutation}}\par
Points fixes de $\sigma_0$
% source matrix, lines 14-20
\[
\deleted{\unsure{\overline{\sigma}_0}=}\quad\sigma_0=
\begin{array}{r@{}l}
&\begin{array}{ccccccc}\underline{1}&\underline{2}&&&&\underline{n+1}&\underline{n+2}\end{array}\\[-.5ex]
\begin{array}{r}1\\2\\\\\\\\n+1\\n+2\end{array}&
\left(\begin{array}{ccccccc}
u_{01}-1&v_{01}-1&v_{01}-1&v_{01}-1&&v_{01}-1&1\\
0&v_{02}&\deleted{\unsure{u_{01}}}0&0&&0&0\\
\deleted{\unsure{1}}0&0&v_{02}&0&&0&0\\
0&\unsure{0}&0&v_{02}&&0&0\\
\hline
0&0&0&0&&\deleted{\unsure{1}}0&0\\
0&0&0&0&&v_{02}&0\\
0&0&0&0&&0&1
\end{array}\right)
\end{array}\tag*{(12)}
\]
% source line 21
\[
\sigma_0(x_1,\ldots,x_{n+1})=\bigl(1+(u_{01}-1)x_1+(v_{01}-1)\overset{\substack{x_2+\cdots+x_{n+1}\\\Vert}}{\Sigma_0},v_{02}x_2,\ldots,v_{02}x_{n+1}\bigr)\tag*{(13)}
\]
% source lines 22-23
\[
H_0:\left\{\begin{aligned}
&1+(u_{01}-2)x_1+(v_{01}-1)\Sigma_1=0\\
&(v_{02}-1)x_2=\cdots=(v_{02}-1)x_{n+1}=0
\end{aligned}\right.\tag*{(14)}
\]
% source lines 24-28
C'est une réflexion projective ssi les dernières relations sont conséquence de la première \unsure{homogénéisée}, ou encore ssi $\boxed{v_{02}=1}$. Il reste pour $\sigma_0$ les paramètres libres $u_{01},v_{01}$, soit $\lambda_0$ et $\mu_0$
% source line 29
\[
\sigma_0(x_1,\ldots,x_{n+1})=\bigl(1+(\lambda_0-1)x_1+(\mu_0-1)\overset{\substack{x_2+\cdots+x_{n+1}\\\Vert}}{\Sigma_0},x_2,\ldots,x_{n+1}\bigr)\tag*{(15)}
\]
